\clearpage
\section*{Preregistered harder split at 3-mer Jaccard 0.40}
\label{sec:harder-split}
\subsection*{Reason for the follow-up}
The earlier cluster-holdout at threshold 0.60 left many near-singleton clusters. The paper's preceding study20 metric thus could not, by itself, prove strong family separation. The dated preregistration addendum (\texttt{docs/PREREG\_ADDENDUM3*}) was committed before the harder-split outcome: use single-linkage 3-mer Jaccard clustering at threshold 0.40, keep whole clusters in one of ten seeded folds, retain the study19 ensemble, and compare the mean MCC to two fixed gates. A mean MCC of at least 0.799 passes the published AMP Scanner Feb2020 point-value hard gate; 0.8328 passes the stronger study19 random-split point-value gate. The neighboring thresholds 0.35 and 0.45 were declared secondary sensitivity arms, not alternatives chosen after observing this primary result. They are pending and are not scored in this paper update.
\subsection*{Clustering diagnostics and source}
The committed \texttt{results/study21\_clusters\_t0.4.json} records 3,721 clusters from the 4,042 records, with 3,536 singleton clusters and a largest cluster of 12. That is still a strongly singleton-dominated partition: a lower Jaccard threshold groups some near matches but does not remove all homology, train-source confounding, or assay bias. The ten held-out fold sizes are 404 or 405 records. These diagnostics are not a proof of independent protein families.
\subsection*{Full primary-arm result}
\texttt{results/study21\_hard\_split\_t0.4.json}, committed on main as study21's final primary arm (8413a88), records all ten per-fold outcomes and mean metrics. Table~\ref{tab:hardfolds} transcribes those held-out results without rerunning or replacing the locked protocol.
\begin{center}\small\begin{tabular}{rrrrrrr}\hline
Fold & n & MCC & Accuracy & AUROC & Sensitivity & Specificity\\\hline
1 & 405 & 0.8369 & 0.9185 & 0.9607 & 0.9184 & 0.9187\\
2 & 405 & 0.8726 & 0.9358 & 0.9747 & 0.9159 & 0.9581\\
3 & 404 & 0.8256 & 0.9134 & 0.9700 & 0.9178 & 0.9081\\
4 & 404 & 0.8372 & 0.9183 & 0.9706 & 0.9479 & 0.8860\\
5 & 404 & 0.8471 & 0.9233 & 0.9716 & 0.9082 & 0.9391\\
6 & 404 & 0.8010 & 0.9010 & 0.9537 & 0.8462 & 0.9459\\
7 & 404 & 0.8479 & 0.9233 & 0.9822 & 0.9482 & 0.9005\\
8 & 404 & 0.8158 & 0.9059 & 0.9614 & 0.8557 & 0.9557\\
9 & 404 & 0.8320 & 0.9158 & 0.9523 & 0.8649 & 0.9589\\
10 & 404 & 0.8214 & 0.9109 & 0.9679 & 0.9108 & 0.9110\\
\hline\end{tabular}\refstepcounter{table}\label{tab:hardfolds}\end{center}
\begin{center}\begin{tabular}{lrrr}\hline
Metric & Tenfold mean & Fold SD & One-sided lower 95\%\\\hline
MCC & 0.8337 & 0.0198 & 0.8223\\
Accuracy & 0.9166 & 0.0098 & 0.9109\\
AUROC & 0.9665 & 0.0094 & 0.9611\\
Sensitivity & 0.9034 & 0.0360 & 0.8825\\
Specificity & 0.9282 & 0.0266 & 0.9128\\\hline
\end{tabular}\end{center}
The recorded gate booleans are true for both mean-MCC comparisons. The stronger numerical margin is only 0.0009 above the 0.8328 target at the precision recorded; fold variability is far larger. It would be wrong to describe this as an unqualified universal state-of-the-art result or an independent significance test against the published method. The threshold 0.40 arm is an in-dataset challenge with a predeclared split, not an external cohort.
\subsection*{Branch-model diagnostic, not a separate benchmark}
The committed mean branch AUROCs are CNN7 0.9573, CNN27 0.9598, GNN 0.8876, random forest 0.9641 and histogram gradient boosting 0.9675. These branches are members of the same ensemble and fold construction, not five independent experiments. Differences cannot be promoted into claims of mechanism or clinical utility without a nested comparison. The GNN is not the strongest branch here; the ensemble's advantage depends on all branch weight choices already fixed in the protocol.
\subsection*{Limitations and next locked checks}
This source reports ten folds of one dataset. The 3-mer single-linkage rule can fail to group distant homologs; input annotations are subject to the review-status/source confound already described in this manuscript. A harder within-dataset split does not validate the ten candidate sequences from the distinct discovery screen. The t=0.35 and t=0.45 secondary sensitivity arms remain [PENDING: committed final fold results]. A truly independent AMP cohort, prospectively fixed before model fitting, and subsequent MIC and hemolysis assays remain [PENDING: external and wet-lab validation]. Neither a candidate score nor an in-silico novelty tier is an antibiotic discovery.
